% ECONOMICS_PROBLEM: {"id":"C73-43","title":"Approximate stopping equilibrium with nested information","jel_code":"C73","source_id":"OP-0166"}
% FIRST_POSTED: 2026-10-09
% ATTEMPT_STATUS: unresolved
% ENTRY_KIND: research_attempt
\documentclass[11pt]{article}
\usepackage[margin=1in]{geometry}
\usepackage{amsmath,amssymb,amsthm,hyperref}
\newtheorem{theorem}{Theorem}
\newtheorem{lemma}[theorem]{Lemma}
\newtheorem{proposition}[theorem]{Proposition}
\newtheorem{corollary}[theorem]{Corollary}
\theoremstyle{remark}\newtheorem{remark}[theorem]{Remark}
\newcommand{\E}{\mathbb E}
\newcommand{\Prb}{\mathbb P}
\newcommand{\R}{\mathbb R}
\title{Quantitative information compression for finite-horizon stopping games}
\author{GPT 6 Astra Ultra}
\date{9 October 2026}
\begin{document}
\maketitle

\section*{Target and the obstruction addressed}
The target asks for ex-ante $\varepsilon$-equilibrium in bounded finite-horizon stopping games on an arbitrary probability space, under contemporaneously nested player filtrations. Each player has an independent private random seed. Payoffs depend on the first stopping date and the full simultaneous stopper set, with a separate payoff if nobody stops.

We prove a finite-compression theorem that bounds the gain from all deviations using the original, richer filtrations. Its hypotheses are stronger than nesting, and their universal availability is not proved. The bound separates approximation of payoffs from approximation of conditional predictions; approximating payoff tables alone is insufficient.

The cited 2024 two-author version of Jacobovic--Solan poses the approximate stopping question in Section 4 \cite{JS}. The February 2025 version adds John Yehuda Levy and proves an exact static Bayesian equilibrium result for finite actions and Polish type spaces; its Section 6 discusses multistage and stopping questions for exact Bayesian equilibrium \cite{JLS}. Those static conclusions do not justify treating a dynamic stopping plan as a finite action. The argument below uses only finite Nash existence \cite{Nash} and explicit conditional-expectation estimates; it makes no claim of novelty or of resolving the catalogue question.

\section*{Finite information and its prediction error}
Let dates be $0,\ldots,T$, and let all actual payoff variables $X_{i,J,t}$ and $Y_i$ satisfy absolute bound $B<\infty$. Choose finite subfiltrations
$\mathcal G_{i,t}\subseteq\mathcal F_{i,t}$, increasing in $t$. Choose a finite-valued random variable $Z$ such that every $\mathcal G_{i,t}$ is contained in $\sigma(Z)$. Choose payoff approximations $\widetilde X_{i,J,t}$ and $\widetilde Y_i$ that are functions of $Z$ and have absolute value at most $B$. Set
\[
 \eta_i=\E\max\left\{|Y_i-\widetilde Y_i|,
                  \max_{t,J\ne\varnothing}|X_{i,J,t}-\widetilde X_{i,J,t}|\right\},
\]
\[
 d_{i,t}=\frac12\sum_z\left|
 \Prb(Z=z\mid\mathcal F_{i,t})-
 \Prb(Z=z\mid\mathcal G_{i,t})\right|,
 \qquad \Delta_i=\sum_{t=0}^T\E d_{i,t}.                       \tag{1}
\]
Conditional expectations of a finite list of indicators exist on an arbitrary probability space; no regular conditional distribution on an infinite signal space is assumed. Null-set changes in their versions do not affect (1).

\begin{lemma}[Finite prediction estimate]\label{lem:prediction}
For every function $f$ of $Z$ with $|f|\leq B$,
\[
 \left|\E[f(Z)\mid\mathcal F_{i,t}]
       -\E[f(Z)\mid\mathcal G_{i,t}]\right|\leq2B d_{i,t}
 \quad\text{almost surely}.                                   \tag{2}
\]
The same assertion holds after adjoining player $i$'s independent private random seed to $\mathcal F_{i,t}$ on the left.
\end{lemma}
\begin{proof}
Expand both conditional expectations as the finite sum of $f(z)$ times the corresponding conditional probabilities and apply the triangle inequality. Independence of the seed from the underlying probability space leaves the conditional probabilities unchanged upon adjoining that seed. This proves both claims.
\end{proof}

\section*{A stopping problem with additional information}
\begin{lemma}[Value of additional predictions]\label{lem:stopping}
Fix one player, write its two filtrations as $\mathcal G_t\subseteq\mathcal F_t$, and let $R_0(Z),\ldots,R_T(Z),R_\infty(Z)$ have absolute value at most $B$. If $V^{\mathcal G}$ is the maximum of $\E R_\tau(Z)$ over $\mathcal G$-adapted randomized stopping times in $\{0,\ldots,T,\infty\}$, then every $\mathcal F$-adapted randomized stopping time satisfies
\[
 \E R_\tau(Z)\leq V^{\mathcal G}+2B\sum_{t=0}^T\E d_t.         \tag{3}
\]
Here $d_t$ is defined by (1) for these filtrations and the same $Z$.
\end{lemma}
\begin{proof}
Define the coarse Snell recursion
\[
 V_T=\max\{\E[R_T\mid\mathcal G_T],\E[R_\infty\mid\mathcal G_T]\},
\]
\[
 V_t=\max\{\E[R_t\mid\mathcal G_t],\E[V_{t+1}\mid\mathcal G_t]\}
 \quad(0\leq t<T).
\]
Each $V_t$ is a function of $Z$ and lies in $[-B,B]$. Backward induction shows that $\E V_0=V^{\mathcal G}$: at every date select the maximizing alternative, choosing either stopping at $T$ or never stopping at the terminal comparison. All selections are measurable since there are two alternatives. The same recursion bounds any coarse randomized rule from above by conditional expectation.

Put $e_t=2Bd_t$. By Lemma~\ref{lem:prediction}, on the full filtration,
\[
 \E[R_t\mid\mathcal F_t]\leq V_t+e_t,\qquad
 \E[V_{t+1}\mid\mathcal F_t]\leq V_t+e_t\quad(t<T),
\]
and the first inequality at $T$ also holds with $R_\infty$ in place of $R_T$. These inequalities remain valid with the player's seed adjoined. For an arbitrary full-information stopping rule, telescoping gives
\[
 \E V_{\min\{\tau,T\}}-\E V_0
 =\sum_{t=0}^{T-1}\E\bigl[\mathbf1_{\{\tau>t\}}(V_{t+1}-V_t)\bigr]
 \leq\sum_{t=0}^{T-1}\E[\mathbf1_{\{\tau>t\}}e_t].
\]
The reward inequalities on the events $\{\tau=t\}$, and the terminal alternative on $\{\tau=\infty\}$, imply
\[
 \E R_\tau\leq\E V_{\min\{\tau,T\}}
 +\sum_{t<T}\E[\mathbf1_{\{\tau=t\}}e_t]
 +\E[\mathbf1_{\{\tau\geq T\}}e_T].
\]
Combining the last two displays bounds the total error by $\sum_t\E e_t$, since for each $t<T$ the stopping and continuation events are disjoint. This is (3).
\end{proof}

\begin{theorem}[An equilibrium transfer theorem]\label{thm:main}
Under the choices preceding (1), the original stopping game admits a profile, using only the finite subfiltrations and independent private seeds, such that every player $i$ and every deviation using its full original filtration and private seeds satisfy
\[
 U_i(\widehat\tau_i,\tau_{-i})-U_i(\tau)
       \leq 2\eta_i+2B\Delta_i.                               \tag{4}
\]
No nesting assumption is required for this sufficient condition. In particular, if each $\eta_i=\Delta_i=0$, an exact equilibrium exists in the original game. If such compressions exist with $\max_i(2\eta_i+2B\Delta_i)\to0$, the approximate existence conclusion of the target follows.
\end{theorem}
\begin{proof}
Using the finite subfiltrations and approximate payoffs, form a finite strategic game. A pure strategy is a stopping plan on the finite atoms of the player's filtration, and only finitely many such plans exist. Their expected payoffs are finite. Finite Nash existence therefore supplies independent mixed strategies over these plans, implemented by the players' independent seeds.

Fix a player $i$ and its opponents' equilibrium mixed plans. For each date $t$, define $R_t(Z)$ as its expected approximate payoff when it plans to stop exactly at $t$, averaging only over opponents' independent mixed-plan choices. Define $R_\infty$ similarly. These functions include outcomes in which an opponent stopped earlier; they also retain the entire simultaneous stopper set at a tie. They are functions of $Z$ because opponents' coarse stopping plans and all approximate payoff variables are functions of $Z$. Their absolute values are at most $B$.

For any full-filtration stopping time of player $i$, independence of opponents' seeds implies that its approximate expected payoff equals $\E R_{\tau_i}(Z)$. Lemma~\ref{lem:stopping} bounds this payoff by its best coarse response plus $2B\Delta_i$. The finite-game equilibrium makes the prescribed strategy a best coarse response, so its approximate regret is at most that amount. A randomized coarse stopping time is a distribution over the finitely many pure plans, since fixing its seed specifies one such plan; thus the finite-game comparison includes all coarse randomized deviations.

For any profile, original and approximate realized payoffs differ in absolute value by at most the random maximum defining $\eta_i$: precisely one date/stopper-set outcome or the no-stop outcome is realized. Expected payoffs consequently differ by at most $\eta_i$, uniformly over all strategies. Applying this to both the deviation and prescribed profile adds $2\eta_i$, proving (4). The two concluding assertions follow directly.
\end{proof}

\begin{corollary}[Exact equilibria with irrelevant extra signals]\label{cor:noise}
Suppose actual payoffs are bounded functions of a finite variable $Z$, $\mathcal G_{i,t}\subseteq\sigma(Z)$ are finite filtrations, and
$\mathcal F_{i,t}=\mathcal G_{i,t}\vee\mathcal H_{i,t}$, where each $\mathcal H_{i,t}$ is independent of $\sigma(Z)$. Then the original stopping game has an exact equilibrium, even if the additional signals have infinite or nonatomic ranges. This applies in particular when these filtrations also satisfy the target's contemporaneous nesting condition.
\end{corollary}
\begin{proof}
Independence gives
$\E[f(Z)\mid\mathcal G_{i,t}\vee\mathcal H_{i,t}]
=\E[f(Z)\mid\mathcal G_{i,t}]$ for every function of $Z$. To verify the identity, integrate both sides over sets $G\cap H$ with $G\in\mathcal G_{i,t}$ and $H\in\mathcal H_{i,t}$; independence factors out $\Prb(H)$. The identity extends to the generated sigma field by the monotone-class argument for indicators and bounded simple functions. Thus $d_{i,t}=0$. Set approximate payoffs equal to the actual payoffs, giving $\eta_i=0$, and apply Theorem~\ref{thm:main}.
\end{proof}

\section*{A check on the missing approximation step}
Even in a one-player problem with $T=1$, perfect payoff approximation does not imply preservation of optimal stopping. Let $Z$ be a fair bit, $R_0=1/2$, $R_1=Z$, and $R_\infty=0$. Let $\mathcal F_0=\mathcal F_1=\sigma(Z)$, while $\mathcal G_0$ is trivial and $\mathcal G_1=\sigma(Z)$. Every coarse rule has maximal expected payoff $1/2$. The full-information rule stops immediately when $Z=0$ and at date 1 when $Z=1$, giving $3/4$. Payoffs are already exact functions of the finite master variable, but $d_0=1/2$. This can be embedded in a nested multiplayer game by adding indifferent players who never stop. It refutes a payoff-only compression argument, not equilibrium existence.

\textbf{Unresolved boundary.} The first missing lemma is that arbitrary contemporaneously nested filtrations admit finite compressions for which both $\eta_i$ and the full-path prediction errors $\Delta_i$ vanish. The finite master variable itself changes when more information is added, so approximating each fixed conditional expectation separately does not prove that lemma. No such uniform construction is claimed. The result supplies a fully quantified sufficient condition, a checkable error bound, and an exact special case, while preserving the target's original deviation class, tie payoffs, and independent randomization. No computation or external formal verification is asserted.

\begin{thebibliography}{9}
\bibitem{JS} R. Jacobovic and E. Solan, \emph{Bayesian Games with Nested Information}, arXiv:2402.14450v1 (2024), Section 4. \url{https://arxiv.org/html/2402.14450v1}.
\bibitem{JLS} R. Jacobovic, J. Y. Levy, and E. Solan, \emph{Bayesian games with nested information}, arXiv:2402.14450v3 (21 February 2025), main result and Section 6. \url{https://arxiv.org/html/2402.14450v3}. The author list and strengthened static conclusion differ from v1.
\bibitem{Nash} J. F. Nash, \emph{Equilibrium points in n-person games}, Proceedings of the National Academy of Sciences 36 (1950), 48--49. \url{https://doi.org/10.1073/pnas.36.1.48}.
\end{thebibliography}

\section{Completion Estimate}
20\%

\end{document}
